%%%%%%%%%%%%%%%%%%%%%%%%%%%%%%%%%%%%%%%%%
% Important note:
% This template requires the resume.cls file to be in the same directory as the
% .tex file. The resume.cls file provides the resume style used for structuring the
% document.
%
%%%%%%%%%%%%%%%%%%%%%%%%%%%%%%%%%%%%%%%%%

%----------------------------------------------------------------------------------------
%	PACKAGES AND OTHER DOCUMENT CONFIGURATIONS
%----------------------------------------------------------------------------------------

\documentclass{resume} % Use the custom resume.cls style

\usepackage{enumitem}
\usepackage{hyperref}
\usepackage[left=0.5in,top=0.6in,right=0.5in,bottom=0.6in]{geometry} % Document margins
\newcommand{\tab}[1]{\hspace{.2667\textwidth}\rlap{#1}}
\newcommand{\itab}[1]{\hspace{0em}\rlap{#1}}
\name{Jasper Geer} % Your name
\address{
  \href{mailto:jasper.geer@ubc.ca}{\textbf{jasper.geer@ubc.ca}}
  \\
  \href{https://www.jaspergeer.com}{\textbf{www.jaspergeer.com}}}
\usepackage{enumitem}
\setlist[itemize]{noitemsep}

\hypersetup{
    colorlinks=true,
    filecolor=magenta,      
    urlcolor=cyan,
    pdfpagemode=FullScreen,
    }

\begin{document}

%----------------------------------------------------------------------------------------
%	EDUCATION SECTION
%----------------------------------------------------------------------------------------

\begin{rSection}{Education}
%--copy and paste this region  if you need more--
{\bf University of British Columbia} \hfill {\em 2024-Present}
\\ {\it Vancouver, British Columbia, Canada}
\\ Computer Science, PhD
\\ Supervised by Prof. Alexander J. Summers

\begin{itemize}
    \item Research in Programming Languages and Formal Verification in the UBC Software Practices Lab.
    \item Interested in Rust verification, separation logic, and type systems.
    \item Member of the Prusti project, building a deductive Rust verifier in collaboration with ETH Zurich.
\end{itemize}

%--copy and paste this region  if you need more--
{\bf Tufts University} \hfill {\em 2020-2024}
\\ {\it Medford, Massachussets, USA}
\\ Computer Science, BS
\\ \textit{summa cum laude}
\end{rSection}
%----------------------------------------------------------------------------------------
%	EXPERIENCE SECTION
%----------------------------------------------------------------------------------------
\begin{rSection}{Research Experiences}
%--copy and paste this region  if you need more--
{\bf Verifying Safe Clients of Interior Mutability in Rust} \hfill {\em August 2025 - Present}\\
with Alexander J. Summers\\
{\it University of British Columbia}
\begin{itemize}
  \item Developed a specification language for types with interior mutability in Rust and an associated client-side verification methodology.
  \item In contrast to existing work, our technique can be applied to existing Rust code using existing types, does not require the user to interact with ghost state or resource algebras, and allows idiomatic runtime assertions to be reused as formal specifications.
  \item Explored formal underpinnings for our technique using Daenerys, an extension of the Iris separation logic framework.
\end{itemize}

{\bf Place Capability Graphs} \hfill {\em January 2025 - 2025}\\
with Alexander J. Summers\\
{\it University of British Columbia}
\begin{itemize}
  \item Developed a MIR-level mutation testing tool for the Rust compiler to evaluate the trustworthiness of Place Capability Graphs, a model of Rust's type checking results.
  \item Co-authored a paper published at OOPSLA 2025.
\end{itemize}

{\bf Type- and Name-Guided Program Synthesis} \hfill {\em September 2023 - September 2025}\\
with Jeffrey Foster \\
{\it Tufts University}
\begin{itemize}
    \item Implemented and formalized a subtyping constraint-guided program synthesis technique for Java.
    \item Developed search heuristics to make a larger class of synthesis problems tractable.
\end{itemize}

{\bf Technical Limitations and Novel Use Cases for Symbolic Execution} \hfill {\em September 2023 - May 2024}\\
with Daniel Votipka\\
{\it Tufts University}

\begin{itemize}
    \item Participated in an extensive review of recent symbolic execution and automated testing literature.
    \item Qualitatively coded rounds of 5-10 research papers and contributed to codebook development.
  \item Co-authored a paper published at SecDev 2026 which received a distinguished paper award.
\end{itemize}
%--copy and paste this region  if you need more--
\end{rSection}

\newpage
\begin{rSection}{Professional Experiences}
{\bf Vehicle Software Intern} \hfill {\em May 2023 - August 2023}\\
\hfill {\it Tesla}
\begin{itemize}
    \item End-to-end feature development in Haskell for an incremental compiler frontend.
    \item Refactored compiler passes into incremental build rules for a monadic build system.
    \item Created embedded domain-specific languages to implement new language server features.
    \item Received offer for full-time conversion.
\end{itemize}
\end{rSection}

\end{rSection}

%--------------------------------------------------------------------------------
%    PUBLICATIONS -----------------------------------------------------------------------------------------------
\begin{rSection}{Publications} 
  {\it 2026}

  {\em SoK: A Modularized Framework for Symbolic Execution and Application for Usable Tool Design}
  \\ James Mattei, Andrew Lin, \textbf{Jasper Geer}, Jie Hu, Moritz Schloegel, Tiffany Bao, Daniel Votipka.
  \\ {\it Secure Development Conference (SecDev)}
  \\ {\bf Distinguished Paper Award}
  \\ \href{https://doi.org/10.1145/3805773.3805996}{DOI}

  {\it 2025}

  {\em Place Capability Graphs: A General-Purpose Model of Rust’s Ownership and Borrowing Guarantees}
  \\ Zachary Grannan, Aurel Bílý, Jonáš Fiala, \textbf{Jasper Geer}, Markus de Medeiros, Peter Müller, Alexander J. Summers.
  \\ {\it Object-oriented Programming, Systems, Languages and Applications (OOPSLA)}
  \\ \href{https://doi.org/10.1145/3763122}{DOI}
\end{rSection}

%--------------------------------------------------------------------------------
%    PREPRINTS -----------------------------------------------------------------------------------------------
\begin{rSection}{Preprints} 
  {\em Mason: Type- and Name-Directed Program Synthesis}
  \\\textbf{Jasper Geer}, Fox Huston, Jeffrey Foster
  \\ {\it arXiv}
  \\ \href{https://doi.org/10.48550/arXiv.2602.16981}{DOI}
\end{rSection}

%--------------------------------------------------------------------------------
%    AWARDS
%-----------------------------------------------------------------------------------------------
\begin{rSection}{Awards} 
{\bf \href{https://www.grad.ubc.ca/awards/four-year-doctoral-fellowship-4yf}{Four Year Doctoral Fellowship}}{, UBC (University-Wide)} \hfill{\em 2026 - 2029}
\begin{itemize}
  \item \$24,000 CAD stipend + full tuition for 4 years
\end{itemize}

{\bf Community Building Award}{, UBC CS Department} \hfill{\em 2025}


\end{rSection}

%--------------------------------------------------------------------------------
%    SERVICE
%-----------------------------------------------------------------------------------------------
\begin{rSection}{Service} 
{\bf Student Development Committee}{, UBC CS Graduate Student Association}\hfill{\em May 2026 - Present}\\
{\bf Reviewer}{, UBC CS Student Research Conference}\hfill{\em 2025, 2026}\\
{\bf VP Social}{, UBC CS Graduate Student Association}\hfill{\em May 2025 - April 2026}\\
{\bf Tuesday Tea Organizer}{, UBC CS Graduate Student Association}\hfill{\em May 2025 - April 2026}\\
{\bf PL Reading Group Organizer}{, UBC Software Practices Lab}\hfill{\em September 2024 - Present}

\end{rSection}

\newpage

\begin{rSection}{Teaching}
{\bf University of British Columbia, Graduate Teaching Assistant} \hfill {\em September 2024 - Present}
\begin{itemize}
    \item CPSC411, Compilers. {\it Spring 2024.}
    \item CPSC311, Definition of Programming Languages. {\it Fall 2024.}
\end{itemize}

{\bf Tufts University, Teaching Fellow} \hfill {\em January 2024 - May 2024}
\begin{itemize}
    \item CS170, Computation Theory. {\it Spring 2024.}
\end{itemize}

{\bf Tufts University, Course Assistant} \hfill {\em September 2022 - December 2023}
\begin{itemize}
    \item CS170, Computation Theory. {\it Fall 2022, Spring 2023, Fall 2023.}
\end{itemize}

{\bf Coding With Kids} \hfill {\em May 2022 - September 2022}
\begin{itemize}
    \item Taught week-long programming classes for primary and secondary school students.
\end{itemize}
    
% %--------------------------------------------------------------------------------
% %    PROJECTS
% %-----------------------------------------------------------------------------------------------
% \begin{rSection}{Projects}
% %--copy and paste this region  if you need more--
% {\bf miniVerifier}
% \begin{itemize}
%     \item Symbolic-execution based verifier for a small programming language with Hoare-logic specifications.
%     \item Written in Haskell and produces SMT-LIBv2 scripts.
% \end{itemize}
% {\bf Compost}
% \begin{itemize}
%     \item LLVM frontend for a statically-typed functional programming language.
%     \item Designed an affine type system to enforce memory safety without runtime garbage collection.
%     \item Began as a personal summer project, completed as a semester-long group project in a compilers class.
%     \item Written in OCaml.
% \end{itemize}
% % {\bf Odin (Contributor)}\\
% % - Contributed to an open source compiler for the Odin programming language \\
% % - Addressed issues in the type checker and code generator \\
% {\bf tinyValidator}
% \begin{itemize}
%     \item Translation validator for a C-subset language by means of symbolic execution.
%     \item Devised a big-step operational semantics to describe the execution of programs with symbolic inputs.
%     \item Written in Haskell and uses the Z3 SMT solver.
% \end{itemize}
% % {\bf jumcc} \\
% % - Compiler from C subset targeting small virtual assembly language \\
% % - Pointers, arrays, functions, structured control flow, and header files \\
% % - Written in Haskell
% %--copy and paste this region  if you need more--
% \end{rSection}
%----------------------------------------------------------------------------------------
%	SKILLS SECTION
%----------------------------------------------------------------------------------------
\begin{rSection}{Programming Background}
Professional/research experience with Rust, Scala, and Haskell.\\
Some experience with OCaml, Standard ML, C, Java, and Rocq.
\end{rSection}

\end{document}----------------------------

